\section{String Theory and Loop Quantum Gravity as HC Sectors}
The decomposition induced by the Heisenberg compensator leads naturally to symmetric and antisymmetric sectors.

\paragraph*{Theorem 9.1 (String--LQG duality via the Heisenberg compensator).}
The Hardy space $H_+$ decomposes under the Heisenberg compensator into two sectors,
\[
H_+ = H_+^{\,+} \oplus H_+^{\,-},
\]
with
\[
C(f)=\tfrac12(f+Jf) \in H_+^{\,+}, \qquad A(f)=\tfrac12(f-Jf) \in H_+^{\,-}.
\]
These sectors are interpreted as follows:
\begin{center}
\begin{tabular}{lll}
\toprule
Sector & Physical theory & Key property \\
\midrule
$H_+^{\,+}$ & String theory & bosonic, $J$-symmetric, continuous \\
$H_+^{\,-}$ & Loop quantum gravity & fermionic, $J$-antisymmetric, discrete \\
\bottomrule
\end{tabular}
\end{center}
Both sectors are unitarily equivalent by Stone--von Neumann~[1]; they are two representations of the same Weyl algebra $W(A_1)$ seen from different frames.

\paragraph*{Identification.}
In the string sector $H_+^{\,+}$, the state $C(f)$ satisfies $J(C(f))=C(f)$ and is therefore $J$-invariant. This mirrors the invariance of the worldsheet under orientation reversal $\sigma \mapsto -\sigma$. In the loop-quantum-gravity sector $H_+^{\,-}$, the state $A(f)$ satisfies $J(A(f))=-A(f)$ and changes sign under $J$, matching the spinorial sign change of oriented holonomies. The discrete LQG area spectrum
\[
A_j = 8\pi \gamma_{\mathrm{BI}} \, \ell_P^2 \sqrt{j(j+1)}, \qquad j \in \tfrac12\mathbb{N},
\]
is consistent with the antisymmetric, Pauli-constrained structure established earlier.

\paragraph*{Proposition 9.2 (Shared vacuum and Planck floor).}
Both sectors share the same vacuum and the same minimal scale:
\begin{itemize}
  \item Vacuum: $z_-=0$ (the photon/Sphaera) is the ground state of both $H_+^{\,+}$ and $H_+^{\,-}$.
  \item Planck floor: $\ell_P = |z_-|_{\min}$ is the string length in $H_+^{\,+}$ and the fundamental area scale in $H_+^{\,-}$.
\end{itemize}
String theory and loop quantum gravity are therefore not competing theories of quantum gravity, but bosonic and fermionic projections of the same object under the Heisenberg compensator.

\paragraph*{Theorem 9.3 (Triolic structure of $\gamma_{\mathrm{BI}}$).}
The Barbero--Immirzi parameter $\gamma_{\mathrm{BI}}$ and the string-sector candidate
\[
\gamma_{\mathrm{string}} = \frac{\ln 2}{\pi\sqrt{3}}
\]
are interpreted as two projections of the same triolic invariant:
\[
\gamma_{\mathrm{string}} + \gamma_{\mathrm{LQG}} + \gamma_3 = 0,
\qquad
\gamma_3 = -\bigl(\gamma_{\mathrm{string}} + \gamma_{\mathrm{LQG}}\bigr) \approx -0.3649.
\]
The frame-invariant quantity is the triolic norm
\[
\tau_\gamma := \sqrt{\frac{\gamma_{\mathrm{string}}^2 + \gamma_{\mathrm{LQG}}^2 + \gamma_3^2}{3}} \approx 0.2619.
\]
A direct computation gives
\[
\gamma_{\mathrm{string}} \approx 0.1274, \qquad \gamma_{\mathrm{LQG}} \approx 0.2375, \qquad \gamma_3 \approx -0.3649,
\]
with vanishing sum.

\paragraph*{Remark (Physical interpretation of the three sectors).}
The three values of $\gamma$ correspond to the three sectors of the biquaternion ground state:
\begin{itemize}
  \item $\gamma_{\mathrm{string}} > 0$: bosonic, the symmetric $C$-sector;
  \item $\gamma_{\mathrm{LQG}} > 0$: fermionic, the antisymmetric $A$-sector;
  \item $\gamma_3 < 0$: the third (tachyonic) sector, projected away by the cascade operator $\hat K$.
\end{itemize}
The negative sign of $\gamma_3$ is not pathological; it is the triolic complement required for the vanishing sum.

\paragraph*{Theorem 9.4 (Wick rotation as splitting mechanism).}
Under the imaginary proper-time transformation $t \mapsto i\tau_E$, the biquaternion time unit transforms as
\[
it \mapsto i(i\tau_E) = -\tau_E \in \mathbb{R}_{<0}.
\]
This identifies the imaginary-time component with the unstable tachyonic sector $\gamma_3<0$. Consequently:
\begin{enumerate}
  \item the balance condition $\gamma_{\mathrm{string}}+\gamma_{\mathrm{LQG}}+\gamma_3=0$ is the stability condition under imaginary proper time;
  \item the positive sectors $\gamma_{\mathrm{string}}, \gamma_{\mathrm{LQG}}>0$ are stable under Wick rotation, whereas the negative sector is projected away by $\hat K$;
  \item the Minkowski norm $|q_n|^2=\tau_n^2-|z|^2$ becomes Euclidean,
  \[
  |q_E|^2 = \tau_E^2 + |z|^2 \ge 0,
  \]
  and the system is unconditionally stable in Euclidean signature.
\end{enumerate}

\paragraph*{Theorem 9.5 (Observer frame offset).}
Every matter observer carries a nonzero contravariant component $z_- \neq 0$. All measured physical constants---including $\hbar$, $G$, $c$, and $\gamma_{\mathrm{BI}}$---are projections of frame-invariant values onto the observer's sector. If $\gamma_0$ denotes the frame-invariant value measured in the photon frame $z_-=0$, and if $\beta=2z_-$ is the observer's boost parameter, then
\[
\gamma_{\mathrm{observed}} = \gamma_0 \cdot \gamma_{\mathrm{Lorentz}}(\beta).
\]
In particular,
\[
\gamma_{\mathrm{BI}} \approx 0.2375
\]
is interpreted as the Lorentz-transformed value of $\gamma_{\mathrm{string}}$, seen from a matter frame with $\beta \approx 0.844$, equivalently $z_- \approx 0.422$.

\paragraph*{Remark (Planck length).}
The Planck length is likewise frame-dependent:
\[
\ell_P^{\mathrm{matter}} = \ell_P^{\mathrm{photon}} \cdot \gamma_{\mathrm{Lorentz}}(\beta) \approx 1.864\, \ell_P^{\mathrm{photon}}.
\]
Thus measured constants are sector projections rather than absolute invariants.

\paragraph*{Theorem 9.6 (Symmetric matter--tachyon frame).}
Matter and tachyon are positioned symmetrically around the photon frame $z_-=0$:
\[
\beta_{\mathrm{matter}} = +\beta_{\mathrm{sym}}, \qquad \beta_{\mathrm{tachyon}} = -\beta_{\mathrm{sym}},
\]
where
\[
\beta_{\mathrm{sym}} = \sqrt{\frac{\gamma_{\mathrm{BI}}-\gamma_{\mathrm{string}}}{\gamma_{\mathrm{BI}}+\gamma_{\mathrm{string}}}} \approx 0.5493.
\]
The Lorentz factor between the matter and tachyon frames is then
\[
\gamma_{\mathrm{matter}\leftrightarrow\mathrm{tachyon}} = \frac{1+\beta_{\mathrm{sym}}^2}{1-\beta_{\mathrm{sym}}^2} = \frac{\gamma_{\mathrm{BI}}}{\gamma_{\mathrm{string}}} \approx 1.864.
\]
Thus the sign-structure interpretation and the frame-dependent interpretation of $\gamma_{\mathrm{BI}}$ are equivalent descriptions of the same frame change.


